\chapter{Краус-мост в родительском действии и проверка гессиана}
\label{ch:v8-kraus-bridge-parent-action-hessian-gate}

\section{Два разных вопроса}

Предыдущая глава доказала существование калибровочно-ковариантного канала
$\mathcal L_{q\ell}$ на алгебре наблюдаемых. Теперь необходимо проверить
более сильное утверждение: возбуждает ли уже принятое классическое действие
Тома VII этот канал в своём вакууме.

Пусть $z_a$ --- вещественные координаты двенадцати межсекторных стрелок, а
$\mathcal L_a=-\operatorname{ad}(D_a)^2/2$. Для нормированного центрального
контраста $Q$ естественная полевая форма Дирихле равна
\begin{equation}
 E_{q\ell}(z)
 =-\langle Q,\mathcal L_zQ\rangle,
 \qquad
 \mathcal L_z=\sum_a z_a^2\mathcal L_a.
 \label{eq:v8-field-dependent-kraus-dirichlet}
\end{equation}
Вычисление даёт одинаковый коэффициент для каждого направления:
\begin{equation}
 E_{q\ell}(z)=\frac7{36}\sum_{a=1}^{12}z_a^2,
 \qquad
 \Hess E_{q\ell}=\frac7{18}I_{12}.
 \label{eq:v8-cross-dirichlet-hessian}
\end{equation}
На восьми контрольных направлениях, не пересекающих два центральных
сектора, коэффициент равен нулю.

\section{Совместимость с запуском Тома VII}

Форма \eqref{eq:v8-cross-dirichlet-hessian} положительна. Поэтому её можно
добавить к родительскому функционалу с любым неотрицательным весом, не
разрушая найденный в Томе VII локальный переход. Машинный аудит для весов
от $10^{-6}$ до $10^6$ сохраняет
\begin{equation}
 (n_-,n_0,n_+)_{\rm origin}=(7,0,20),
 \qquad
 (n_-,n_0,n_+)_{\rm vac}=(0,0,27).
\end{equation}
Но это уточняет физический смысл результата: добавка увеличивает массы
двенадцати тяжёлых межсекторных направлений. Она не делает их
неустойчивыми и не создаёт ненулевой классический фон.

\section{Почему канал в вакууме выключен}

В целевом вакууме Тома VII все двенадцать межсекторных стрелок координат равны
нулю. Поэтому
\begin{equation}
 z_a=0\quad\Longrightarrow\quad
 \mathcal L_z=0,
 \qquad
 \Gamma_{q\ell}^{\rm tree}=0.
 \label{eq:v8-tree-kraus-rate-zero}
\end{equation}
Ненулевой гессиан означает только стоимость флуктуации. Он не равен
ненулевой интенсивности флуктуаций.

Чтобы получить работающий канал, требуется положительная ковариация
$C_{ab}=\langle z_az_b\rangle$. Калибровочная симметрия оставляет скалярный вес на
каждом полном мультиплете, но не отождествляет две семьи $QLYR$ и $XLdR$.
Поэтому остаются как минимум два числа $c_Q,c_X$. Центральное ядро при
$c_Q,c_X>0$ одномерно, а скорость равна
\begin{equation}
 \Gamma_{q\ell}(c_Q,c_X)=\frac76(c_Q+c_X).
 \label{eq:v8-kraus-rate-covariance-scale}
\end{equation}
Даже гауссовская подстановка
$C_T=T(H_{\rm heavy})^{-1}$ сохраняет свободную общую силу флуктуаций
$T$; при $T=1$ выбранная нормировка даёт $35/96$, но это не физически
выведенное число.

\begin{theorem}[Кинематический мост без классического запуска]
Усреднённый по калибровочной группе межсекторный Краус-мост совместим с качественным
гессианом Тома VII: его положительная форма Дирихле сохраняет семь
неустойчивых корневых и двадцать тяжёлых направлений в начале, а также
строгую устойчивость целевого вакуума. Однако в самом классическом вакууме
межсекторная скорость равна нулю. Ненулевой процесс требует ковариации,
квантового усреднения или физической дилатации, не содержащихся в
классическом одноточечном поле.
\end{theorem}

\begin{nogocriterion}[Запрет подмены массы шумом]
Нельзя объявлять положительную квадратичную массу межсекторной стрелки
готовой Краус-скоростью. До предъявления меры флуктуаций, состояния среды
или внутреннего квантового правила усреднения канал остаётся кинематически
допустимым, но динамически выключенным на древесном вакууме.
\end{nogocriterion}

\begin{protocol}
\begin{enumerate}
 \item полевая форма Дирихле: \textbf{выведена};
 \item её опора: \textbf{12 межсекторных тяжёлых направлений};
 \item коэффициент одного направления: \textbf{$7/36$};
 \item сигнатуры Тома VII при положительном весе: \textbf{сохранены};
 \item запуск межсекторных стрелок конденсата: \textbf{нет};
 \item древесная скорость в вакууме: \textbf{ноль};
 \item ковариация двух семейств: \textbf{имеет как минимум два свободных веса};
 \item единственная физическая скорость: \textbf{не выведена};
 \item следующий вопрос: \textbf{внутреннее происхождение ковариации}.
\end{enumerate}
\end{protocol}

Машинный сертификат сохранён в
\path{s2t_v8_kraus_bridge_parent_action_hessian_gate_results.json}.

\begin{thebibliography}{9}
\bibitem{v8-parent-davies}
E.~B.~Davies,
\emph{Markovian master equations},
Communications in Mathematical Physics \textbf{39} (1974), 91--110.
\bibitem{v8-parent-covariant-stinespring}
D.~Verdon,
\emph{A covariant Stinespring theorem},
Journal of Mathematical Physics \textbf{63} (2022), 091705.
\bibitem{v8-parent-open-systems}
H.-P.~Breuer, F.~Petruccione,
\emph{The Theory of Open Quantum Systems}, Oxford University Press, 2002.
\end{thebibliography}